\section{Related Work}
\label{sec:related}

\paragraph{Multilingual and dialect adaptation.}
Cross-lingual encoders such as multilingual BERT and XLM-R support transfer
from high-resource languages \citep{devlin2019bert,conneau2020xlmr}. Adapter
methods further isolate language-specific parameters while sharing a common
backbone \citep{pfeiffer2020madx}. Byte-level modeling avoids unknown tokens
and represents productive spelling variation directly
\citep{xue2022byt5}. Our method is complementary: it changes neither the
tokenization nor the backbone and instead adds a small, editable memory.

\paragraph{Retrieval-augmented generation.}
Retrieval-augmented language models condition generation on external passages
\citep{lewis2020rag}; nearest-neighbor machine translation interpolates a
model distribution with target tokens retrieved from a datastore
\citep{khandelwal2021knnmt}. These systems typically retrieve for every input.
\method instead makes retrieval conditional on calibrated uncertainty and
tests whether evidence is safe to use.

\paragraph{Responsible language technology.}
Aggregate benchmark scores can conceal uneven quality across language
varieties and social groups \citep{blodgett2020language}. We therefore report
per-variety outcomes, separate named-entity preservation from surface
accuracy, and document the assumptions behind the illustrative benchmark.

% -----------------------------------------------------------------------------
